\cotefile{134-2}
\scanpage{146}\mspage{1}
\hfill Villecun 17/19 \unsure{Juillet|juillet}\unsure{.}\\
\hfill 1975\\
Cher Larry,\\

Tout d'abord félicitations pour la naissance de ta fille, qui heureusement ne s'est pas laissée impressionner par les importantes cogitations mathématiques de son papa -- au risque de les perturber quelque peu.

Au fait, quelle est ta situation en France : es\deleted{\unsure{t}}-tu français ou américain, prof. associé ou titulaire ? Si cela pouvait s'arranger, \unsure{serait|serais} tu intéressé par une position : Montpellier ? Cela te dit-il quelque chose de venir à l'occasion à Montpellier -- ou à Villecun -- pour discuter de tes favoris avec moi -- pour le moment il n'y a guère que moi ici pour ne pas nager entièrement (seulement partiellement) dans ce que tu fais.

J'ai été content de finir par recevoir un écho à ma longue \unsure{èpitre|épitre} -- et même un début d'approche constructive pour \deleted{\illegible{1}} une théorie du type que je prévois. La construction que tu proposes de la notion de $n$-catégorie pas stricte, et du foncteur nerf, a certes le mérite \inserted{d'exister et} d'être une premièr approche précise, mais par ailleurs prête \deleted{\unsure{naturellement|évidemment}} \inserted{bien sûr}, le flanc à des critiques évidentes : d'être très technique, peu intuitive (alors \unsure{qu'au|où au} niveau des $1$-cat. et
% archive sheet number: 146
